%%
%% 
%% Style file for MECH MSc thesis. 

%%

\documentclass[11pt,a4paper]{article}

\usepackage{graphicx}
\usepackage{amsmath}
\usepackage{hyperref}
\usepackage{enumitem}
\usepackage{geometry}
\usepackage{fancyhdr}

\usepackage{sectsty}
\usepackage[T1]{fontenc}
\usepackage{titlesec}
\sectionfont{\centering\MakeUppercase}

\geometry{top=2.54cm, bottom=2.54cm, left=1.91cm, right=1.91cm}

\title{Project Title}
\author{Author Title\\
        MSc Programme\\
        UCL Mechanical Engineering\\
        London WC1E7JE}
        
\begin{document}

\maketitle

\begin{abstract}
This document provides a guide for preparing a thesis paper (on style and structure). You can use this document as a template if you are using Microsoft Word. The LaTeX template is available as well. Instead of using the templates, you can use this document as an instruction set. Titles should be written in uppercase and lowercase letters, not all uppercase. The full name of a student should be stated as an author. In the project title, all variables should appear in lightface italics; numbers and units will remain bold. An abstract should be a single paragraph (up to 300 words). For an abstract to be effective, it must be an accurate stand-alone reflection of the contents of the article. It should contain a brief introduction to the work that is reported and the main conclusions. They should not contain displayed mathematical equations, numbered reference citations, or footnotes. Ensure that your MSc Thesis reads well and is grammatically correct.
\end{abstract}

\textbf{Keywords:} Enter keywords or phrases in alphabetical order separated by commas.

\pagestyle{fancy}
\fancyhead[L]{UCL Mechanical Engineering}
\fancyhead[C]{MSc Thesis}
\fancyhead[R]{AY23-24}
\renewcommand{\headrulewidth}{0pt}
\renewcommand{\footrulewidth}{0pt}

\section{Introduction}
This document presents the instructions for writing an MSc Thesis (Version 24.0). It can also be used as a template for Microsoft Word. The LaTeX template is also available on Moodle (or you can use the Overleaf editor and create an account using your UCL credentials). If you use this document as a template, you should check the final PDF version before submission. The page count in the template is an estimate; refer to the prescribed thesis length for your particular programme. Do not adjust line and character spacing to fit your paper to a specific length specified.

\section{The Narrative}
When planning the thesis (paper article), you should focus on its purpose and the audience:
\begin{itemize}
    \item The markers (including the External Examiners) will be most interested in assessing the quality of your work, knowledge demonstrated, and presentation so they can assess your report against the marking rubric.
    \item Future readers may be staff or students at UCL or possibly the wider engineering community who may not have been involved in the course, may not have seen the equipment, and/or may never use the computer program used.
\end{itemize}
The examiners will look for clearly written articles that are structured and easy to follow. There must be clear evidence that the reported work is your own. Make a clear distinction between the author's work and the work of others. The article needs to highlight how it relates to the MSc programme you are enrolled on. For instance, students in the MSc Engineering with Finance should have both a finance component and an engineering component reflected within their thesis. The article should also reflect on the societal impact (sustainability, safety, diversity and inclusion, and ethics perspectives).

\section{The Article Structure}
The article should start with the abstract and the nomenclature at the beginning and finish with the References. While the author will develop a structure suitable to the particular project and narrative, the following are the core elements to include:
\subsection{Introduction}
The introduction should explain the background of the topic and the problem that will be addressed with this project. The technical level at the start should be aimed at readers with sound general engineering backgrounds but not specialist understanding.
\subsection{Literature review}
The literature review should cover the relevant references within the subject matter: theory background, similar work from literature related to the topic or methods that will be used. Once the gap in the literature is identified, the aims and objectives of your project should be clearly stated.
\subsection{Objectives and research/design approach (i.e., methodology)}
These two sections are very important to help the reader understand what you are trying to achieve and how, hence making it easier to follow the coming sections.
\subsection{The student work and results}
The student should think carefully about how to communicate their work well. Not all work needs to be presented to the reader in the main body of the text. The work and results may be presented in both the main sections and appendix (experimental/ analytical/ numerical/ design work). Focus on the most relevant work, while the rest can be given as additional information in Appendices only.
\subsection{Validations, verification, and/or case study}
The quality assessment of the findings is an essential element of your project. Depending on the project, this may be validation, verification, testing, and/or a case study to prove that objectives are adequately achieved. In case these assessments raise concerns (and there is no time for further work), they should be highlighted and further discussed to demonstrate your knowledge and critical thinking.
\subsection{Discussion}
The discussions section should bring all the reported work together, reflecting on the theory and referencing to evidence from the work or literature. The discussion should be aligned with the project objectives, so make sure all are covered. The discussion allows for critical analysis and hypothesising which will provide evidence of the in-depth learning needed for the high band of marks.
\subsection{Conclusion}
A conclusion should be brief. It might elaborate on the importance of the work, confirm the aim is achieved, and/or suggest applications and further work.

\section{Use of Generative Artificial Intelligence (AI)}
The use of AI generative tools for the MSc Individual Project is limited to Category 2: AI tools can be used in an assistive role only and students should be very careful if using them. Discuss with your supervisor beforehand to avoid any potential academic misconduct. These tools should not be used to write or translate the thesis. In some particular projects, AI tools may present an integral part of the thesis; their use should be specified in the approach as one of the applied methods. These may be used only with the supervisor's approval and the Module Coordinator should be informed. Any use of AI sources in the work should be specifically acknowledged, clearly explaining how they have been used and to which extent. The instructions can be found here: How to acknowledge AI sources in your work.

\section{Guidelines For Article Style}
This template is intended as a template only, not a final article paper. The author should ensure before the submission that the final version in PDF format is of a good style. The page layout should be A4 with the top and bottom margins of 2.54 cm and right and left margins of 1.91 cm. The paragraph should have single line spacing and indentation first line 0.25 cm. The font should be Times New Roman with the font size 11 and not be less than 9 for any text in the article (e.g., tables). Use italics for emphasis; do not underline. For more detailed style information, you may find of interest the IEEE Editorial Style Manual for Authors. This contains a formal set of editorial guidelines including:
\begin{itemize}
    \item punctuation;
    \item capitalization;
    \item abbreviations;
    \item section headings;
    \item numbers, equations;
    \item footnotes;
    \item biographies;
    \item some common mistakes;
    \item units of measurement.
\end{itemize}

The title of the article should accurately encompass the presented work and topic. If during the writing of this article it is identified that the title should be changed to be more suitable, confirm with your supervisor and inform the administration team \href{mailto:mecheng.pg-queries@ucl.ac.uk}{mecheng.pg-queries@ucl.ac.uk} (cc-ing your supervisor) with the subject “MSc project title change <student number>”.

\subsection{Abbreviations and Acronyms}
Define abbreviations and acronyms the first time they are used in the text, even after they have already been defined in the abstract. Abbreviations that incorporate periods should not have spaces: write “C.N.R.S.” not “C. N. R. S.” Do not use abbreviations in the title unless they are unavoidable.

\subsection{Math Presentation}
The equations need to be properly written. You can use either the Microsoft Equation Editor or a similar.
\subsubsection{Equations}
Number equations consecutively with equation numbers in parentheses flush with the right margin of the column as in (1). First use the equation editor to create the equation. Then you can select the “Equation” markup style. Press the tab key and write the equation number in parentheses. To make your equations more compact, you may use the solidus ( / ), the exp function, or appropriate exponents. Use parentheses to avoid ambiguities in denominators. Punctuate equations when they are part of a sentence as in

\begin{equation}
B_p + H_2 = 40. \quad (1)
\end{equation}

Be sure that the symbols in your equation have been defined before the equation appears or immediately following. Italicize symbols (T might refer to temperature, but T is the unit tesla). When referring to an equation or formula, use “(1)” not “Eq. (1)” except at the beginning of a sentence: “Equation (1) is ... .”

\subsubsection{Algorithms and codes}
Algorithms and codes should be numbered and include a short title. They are set off from the text with rules above and below the title and after the last line.

\begin{figure*}[h!]
\centering
\includegraphics[width=0.5\textwidth]{figures/Picture1.png}
\end{figure*}

\begin{figure}[h!]
\centering
\includegraphics[width=0.5\textwidth]{figures/figure1.png}
\caption{This is a sample of a figure.}
\end{figure}

\section{Guidelines for Graphics Preparation}
\subsection{Types of Graphics}
The figures should be placed at the top of the page closest to the first mention in the text. Place figure captions below the figures; place table headings above the tables. Do not include captions as part of the figures or put them in “text boxes” linked to the figures. Also, do not place borders around the outside of your figures. When referencing your figures and tables within your article, use the abbreviation “Fig.” even at the beginning of a sentence. Do not abbreviate “Table.” Tables should be numbered with Roman numerals. Figures may include photographs, illustrations, multicolor graphs, and flowcharts and should have well-defined captions (Fig. 1 and 2). Tables and data charts can be in black and white or in color. They should be clear and well-formatted as shown in Table 1. The charts, graphs, and tables, as shown in Figure 3, need to stay within the margins. When choosing the depth of a graphic, allow space for a caption. When preparing your graphics, use Times New Roman if possible for consistency or Symbol. Check that the final version of the graphics (equations, figures, etc.) is not distorted or missing.

\subsection{Multipart Figures}
These are figures compiled of more than one sub-figure presented side-by-side or stacked. They need to have well-defined captions as shown in Fig. 3.

\begin{table}[h!]
\centering
\caption{This is a Sample of a Table Title}
\begin{tabular}{c|c|c|c|c|c}
\textbf{Name} & \textbf{\#1} & \textbf{\#2} & \textbf{\#3} & \textbf{\#4} & \textbf{\#5} \\ \hline
RVF         & 10            & 87            & 55            & 55            & 45            \\ \hline
DUM-PT      & 
33   & 97            & 77            & 68            & 37   \\ \hline
PT-PK       & 55            & 12   & 47   & 91            & 22            \\ \hline
SCPS-CK     & 78            & 10            & 98            & 99            & 88   \\ \hline
LPK-mJJ     & 14            & 88            & 31            & 67   & 82            \\ \hline
Placebo     & 17            & 44            & 70            & 100           & 48   \\ \hline
\end{tabular}
\end{table}

\begin{figure}[h!]
\centering
\includegraphics[width=0.5\textwidth]{figures/figure2.png}
\caption{Sample graph of the diagram}
\end{figure}

\begin{figure}[h!]
\centering
\includegraphics[width=0.9\textwidth]{figures/figure3.png}
\caption{Sample figure (a) Iso-surfaces of Q-criteria showing vortex systems (b) Non-dimensional vorticity components}
\end{figure}

\subsection{Using Labels Within Figures}
\subsubsection{Figure Axis Labels}
Figure axis labels are often a source of confusion. Use words rather than symbols. As an example, write the quantity “Magnetization” or “Magnetization M” not just “M.” Put units in parentheses. Do not label axes only with units. For example, write “Magnetization (A/m)” or “Magnetization (Am$^{-1}$)” not just “A/m.” Do not label axes with a ratio of quantities and units. For example, write “Temperature (K)” not “Temperature/K.” Multipliers can be especially confusing. Write “Magnetization (kA/m)” or “Magnetization ($10^3$ A/m).” Do not write “Magnetization (A/m) × 1000” because the reader would not know whether the top axis label means 16000 A/m or 0.016 A/m. Figure labels should be legible approximately 8- to 10-point type.

\subsubsection{Caption in Multipart Figures and Tables}
Multipart figures should be combined and labeled before final submission. Labels should appear centered below each subfigure in 8-point Times New Roman font in the format of (a) (b) (c).

\subsubsection{File Formats for Graphics}
Format and save your graphics using a suitable graphics processing program that will allow you to create the images as PostScript (PS), Encapsulated PostScript (.EPS), Tagged Image File Format (.TIFF), Portable Document Format (.PDF), JPEG, or Portable Network Graphics (.PNG). These programs can re-size them and adjust the resolution settings. The proper resolution of your figures will depend on the type of figure that should be clear (min 300 dpi).

\chapter{Thesis Submission}
The thesis article should be submitted on Moodle in PDF format by the defined deadline. Any relevant additional files (e.g., video, code, modeling, taken databases, experimental data, numerical output files, etc.) can be submitted on Moodle as well. Please check Moodle submission file limits. In case the files are too large to be submitted on Moodle, contact the administration team to arrange a data transfer link. All data files should be submitted by the deadline as well and named appropriately (candidate number, date, title).

\chapter{Conclusion}
A conclusion should state the main points of the article but not replicate the abstract. It should be short, strong, and to the point.

\chapter{Acknowledgement}
These are typically written as “The author thanks ... .” In most cases a supervisor, mentor, sponsor, external help, etc. In addition, any relevant work aspect not the student's original contribution should be identified here (e.g., use of generative AI tools, lab staff).

\chapter{Appendix}
The appendix can be added if needed. Its role is to support the work presented in the main body of the article. The appendix should be prepared in the same style and quality as the rest of the article. Note the maximum size of the Appendix defined.

\chapter{References}
References should be cited in the text and use the IEEE reference style as shown in this document. Do not use “Ref.” or “reference” except at the beginning of a sentence: “Reference [3] shows ... .” You can use the “References” tab in Word to apply them and create a reference list at the end of the paper or write them manually. In all references, the given name of the author or editor is abbreviated to the initial only and precedes the last name. References may not include all information; please obtain and include relevant information. If there is a URL included with the reference, it can be included at the end of the reference. When using multiple references, each should be numbered with separate brackets. When citing a section in a book, please give the relevant page numbers. In the text, refer simply to the reference number.

The basic format for periodicals:
\begin{itemize}
    \item J. K. Author, “Name of paper,” \textit{Abbrev. Title of Periodical}, vol. x, no. x, pp. xxx-xxx, Abbrev. Month, year, doi: 10.1109.XXX.1234567.
\end{itemize}
Examples:
\begin{itemize}
    \item J. U. Duncombe, “Infrared navigation—Part I: An assessment of feasibility,” \textit{IEEE Trans. Electron Devices}, vol. ED-11, no. 1, pp. 34–39, Jan. 1959, doi: 10.1109/TED.2016.2628402.
    \item P. Kopyt \textit{et al.}, “Electric properties of graphene-based conductive layers from DC up to terahertz range,” \textit{IEEE THz Sci. Technol.}, doi: 10.1109/TTHZ.2016.2544142.
    \item D. Comite and N. Pierdicca, "Decorrelation of the near-specular land scattering in bistatic radar systems," \textit{IEEE Trans. Geosci. Remote Sens.}, early access, doi: 10.1109/TGRS.2021.3072864. (Note: This format is used for articles in early access. The DOI must be included.)
    \item H. V. Habi and H. Messer, "Recurrent neural network for rain estimation using commercial microwave links," \textit{IEEE Trans. Geosci. Remote Sens.}, vol. 59, no. 5, pp. 3672-3681, May 2021. [Online]. Available: \url{https://ieeexplore.ieee.org/document/9153027}
\end{itemize}

The basic format for books:
\begin{itemize}
    \item J. K. Author, “Title of chapter in the book,” in \textit{Title of Published Book}, xth ed. City of Publisher (only U.S. State), Country: Abbrev. of Publisher, year, ch. x, sec. x, pp. xxx–xxx.
\end{itemize}
Examples:
\begin{itemize}
    \item G. O. Young, “Synthetic structure of industrial plastics,” in \textit{Plastics}, 2nd ed., vol. 3, J. Peters, Ed. New York, NY, USA: McGraw-Hill, 1964, pp. 15–64.
    \item W.-K. Chen, \textit{Linear Networks and Systems}. Belmont, CA, USA: Wadsworth, 1993, pp. 123–135.
    \item Philip B. Kurland and Ralph Lerner, eds., \textit{The Founders’ Constitution}. Chicago, IL, USA: Univ. of Chicago Press, 1987, Accessed on: Feb. 28, 2010 [Online]. Available: \url{http://press-pubs.uchicago.edu/founders/}
\end{itemize}

The basic format for handbooks:
\begin{itemize}
    \item Name of Manual/Handbook, x ed., Abbrev. Name of Co., City of Co., Abbrev. State, Country, year, pp. xxx-xxx.
\end{itemize}
Examples:
\begin{itemize}
    \item Motorola Semiconductor Data Manual, Motorola Semiconductor Products Inc., Phoenix, AZ, USA, 1989.
    \item R. J. Hijmans and J. van Etten, “Raster: Geographic analysis and modelling with raster data,” R Package Version 2.0-12, Jan. 12, 2012. [Online]. Available: \url{http://CRAN.R-project.org/package=raster}
\end{itemize}

The basic format for reports:
\begin{itemize}
    \item J. K. Author, “Title of report,” Abbrev. Name of Co., City of Co., Abbrev. State, Country, Rep. xxx, year.
\end{itemize}
Example:
\begin{itemize}
    \item E. E. Reber, R. L. Michell, and C. J. Carter, “Oxygen absorption in the earth’s atmosphere,” Aerospace Corp., Los Angeles, CA, USA, Tech. Rep. TR-0200 (4230-46)-3, Nov. 1988.
\end{itemize}

Basic format for conference proceedings:
\begin{itemize}
    \item J. K. Author, “Title of paper,” in \textit{Abbreviated Name of Conf.}, City of Conf., Abbrev. State (if given), Country, year, pp. xxx-xxx.
\end{itemize}
Examples:
\begin{itemize}
    \item D. B. Payne and J. R. Stern, “Wavelength-switched passively coupled single-mode optical network,” in \textit{Proc. IOOC-ECOC}, Boston, MA, USA, 1985, pp. 585–590.
    \item D. Ebehard and E. Voges, “Digital single sideband detection for interferometric sensors,” presented at the 2nd Int. Conf. Optical Fiber Sensors, Stuttgart, Germany, Jan. 2-5, 1984.
    \item PROCESS Corporation, Boston, MA, USA. Intranets: Internet technologies deployed behind the firewall for corporate productivity. Presented at INET96 Annual Meeting. [Online]. Available: \url{http://home.process.com/Intranets/wp2.htp}
\end{itemize}

Basic format for electronic documents (when available online): 
\begin{itemize}
    \item Issuing Organization. (year, month day). \textit{Title}. [Type of medium]. Available: site/path/file
\end{itemize}
Example:
\begin{itemize}
    \item U.S. House. 102nd Congress, 1st Session. (1991, Jan. 11). H. Con. Res. 1 Sense of the Congress on Approval of Military Action. [Online]. Available: LEXIS Library: GENFED File: BILLS
\end{itemize}

Basic format for theses (M.S.) and dissertations (Ph.D.):
\begin{itemize}
    \item J. K. Author, “Title of thesis,” M.S. thesis, Abbrev. Dept., Abbrev. Univ., City of Univ., Abbrev. State, year.
\end{itemize}
Examples:
\begin{itemize}
    \item J. O. Williams, “Narrow-band analyzer,” Ph.D. dissertation, Dept. Elect. Eng., Harvard Univ., Cambridge, MA, USA, 1993.
    \item N. Kawasaki, “Parametric study of thermal and chemical nonequilibrium nozzle flow,” M.S. thesis, Dept. Electron. Eng., Osaka Univ., Osaka, Japan, 1993.
\end{itemize}

The basic format for the most common types of unpublished references:
\begin{itemize}
    \item J. K. Author, private communication, Abbrev. Month, year.
    \item J. K. Author, “Title of paper,” unpublished.
    \item J. K. Author, “Title of paper,” to be published.
\end{itemize}
Examples:
\begin{itemize}
    \item A. Harrison, private communication, May 1995.
    \item A. Brahms, “Representation error for real numbers in binary computer arithmetic,” IEEE Computer Group Repository Paper R-67-85.
\end{itemize}

The basic format for datasets:
\begin{itemize}
    \item Author, Date Year. “Title of Dataset,” distributed by Publisher/Distributor, \url{http://url.com} (or if DOI is used end with a period)
\end{itemize}
Example:
\begin{itemize}
    \item U.S. Department of Health and Human Services, Aug. 2013, “Treatment Episode Dataset: Discharges (TEDS-D): Concatenated 2006 to 2009,” U.S. Department of Health and Human Services, Substance Abuse and Mental Health Services Administration, Office of Applied Studies, doi: 10.3886/ICPSR30122.v2.
\end{itemize}


\bibliographystyle{IEEEtran}

% \bibliographystyle{ieeetr}
\bibliography{manuscript_references}



\end{document}
